% =====================================================================
% Paper 01 — Pillar-1 conditional programme: BCC ranking + finite-N
%            mass-gap anchor + analytic shell-gap bound
%   (NOT a "rigorous Pillar-1 closure paper" at the present canonical tier)
% Status: [DRAFT] [NEEDS_UPDATE] — Wave 1
% AUDIT-FLAG: AUDIT-2026-05-02-Wave7-Aux-Epoch-Overclaim
%   (Math314-AddC extension covering Wave 1/3 papers Paper-00..05).
%   Hostile-referee audit (2026-05-02) flagged this paper as the
%   STRONGEST over-claim of Wave 1/3: "rigorous uniqueness of the BCC
%   ground state", "establish mass-gap closure (Pillar 1)", "stable
%   minimum", "combination establishes mass-gap closure to tier PROVED
%   CONDITIONAL". Actual content: (i) competitor enumeration, (ii)
%   mean-field Brazovskii ranking (single-shell SMA), (iii) finite-N
%   numerical anchor at single operating point, (iv) analytic shell-gap
%   bound. Recast as CONDITIONAL Pillar-1 PROGRAMME, NOT a closure
%   theorem; title and abstract recast.
% Compile: pdflatex Paper-01.tex (×2 for refs)
% Canonical archive: Math01-v2, Math82-AddF, Math82-AddG, Math82-AddH,
%                    Math314-AddC
% =====================================================================
% --- TECT canonical preamble (see _shared/tect-preamble.tex) ----------
% =====================================================================
% TECT canonical LaTeX preamble (revtex4-2 / single-column / theorem-safe)
% =====================================================================
% Maintainer: Jusang Lee (jtkor@outlook.com)
% Binding from: 2026-05-06
% Purpose: a single source-of-truth preamble for every TECT paper
% (Paper-00 .. Paper-10 + Auxiliary notes). Designed to (a) eliminate
% font-corruption on pdflatex (T1 fontenc + lmodern), (b) make
% theorem-heavy mathematical-physics content render cleanly in a
% single-column layout (PRD class option, NOT PRL two-column), (c)
% provide consistent theorem numbering across all TECT papers via
% amsthm with a shared counter set, (d) make hyperref + cleveref
% cross-references resolve cleanly with the standard two-pass
% pdflatex compile.
%
% Usage in each Paper-NN.tex:
%   % =====================================================================
% TECT canonical LaTeX preamble (revtex4-2 / single-column / theorem-safe)
% =====================================================================
% Maintainer: Jusang Lee (jtkor@outlook.com)
% Binding from: 2026-05-06
% Purpose: a single source-of-truth preamble for every TECT paper
% (Paper-00 .. Paper-10 + Auxiliary notes). Designed to (a) eliminate
% font-corruption on pdflatex (T1 fontenc + lmodern), (b) make
% theorem-heavy mathematical-physics content render cleanly in a
% single-column layout (PRD class option, NOT PRL two-column), (c)
% provide consistent theorem numbering across all TECT papers via
% amsthm with a shared counter set, (d) make hyperref + cleveref
% cross-references resolve cleanly with the standard two-pass
% pdflatex compile.
%
% Usage in each Paper-NN.tex:
%   % =====================================================================
% TECT canonical LaTeX preamble (revtex4-2 / single-column / theorem-safe)
% =====================================================================
% Maintainer: Jusang Lee (jtkor@outlook.com)
% Binding from: 2026-05-06
% Purpose: a single source-of-truth preamble for every TECT paper
% (Paper-00 .. Paper-10 + Auxiliary notes). Designed to (a) eliminate
% font-corruption on pdflatex (T1 fontenc + lmodern), (b) make
% theorem-heavy mathematical-physics content render cleanly in a
% single-column layout (PRD class option, NOT PRL two-column), (c)
% provide consistent theorem numbering across all TECT papers via
% amsthm with a shared counter set, (d) make hyperref + cleveref
% cross-references resolve cleanly with the standard two-pass
% pdflatex compile.
%
% Usage in each Paper-NN.tex:
%   \input{../_shared/tect-preamble.tex}
%   \begin{document}
%   ...
%   \end{document}
%
% Build (every paper): `pdflatex Paper-NN && pdflatex Paper-NN`
% (the second pass resolves \ref{} / \cref{} forward references).
% A bash/PowerShell wrapper that automates this is at
% Docs/papers/papers/_shared/build-paper.sh / build-paper.ps1.
% =====================================================================

% ---- Document class --------------------------------------------------
% PRD class option of revtex4-2 with [reprint,onecolumn] gives the same
% APS heading / by-line / abstract style as PRL but in a wider single
% column that handles long display equations and theorem environments
% without column-break artefacts. nofootinbib keeps citations out of
% footnotes. The 11pt option restores standard font metrics; with T1
% fontenc + lmodern this gives non-bitmap glyphs.
\documentclass[%
  aps,prd,reprint,onecolumn,nofootinbib,11pt,%
  amsmath,amssymb,floatfix,longbibliography%
]{revtex4-2}

% ---- Encoding & fonts ------------------------------------------------
\usepackage[T1]{fontenc}        % proper 8-bit font encoding (no font corruption)
\usepackage[utf8]{inputenc}     % allow UTF-8 source
\usepackage{lmodern}            % scalable Latin Modern (replaces bitmap CM)
\usepackage{textcomp}           % extra text symbols
\usepackage{microtype}          % subtle kerning improvements

% ---- UTF-8 → LaTeX character mappings --------------------------------
% Maps the Unicode characters that occur in TECT body text (legacy from
% the chat-content auto-archival rule, CLAUDE.md §4) to LaTeX-safe
% commands. Without these declarations, T1+inputenc rejects the chars
% with "Unicode character X not set up for use with LaTeX".
% Adding declarations centrally (here) is preferable to per-file edits.
\DeclareUnicodeCharacter{00A7}{\S}              % §  (section sign)
\DeclareUnicodeCharacter{00B2}{\ensuremath{^{2}}}        % ²
\DeclareUnicodeCharacter{00B3}{\ensuremath{^{3}}}        % ³
\DeclareUnicodeCharacter{00B5}{\ensuremath{\mu}}         % µ (micro)
\DeclareUnicodeCharacter{00B7}{\ensuremath{\cdot}}       % · (middle dot)
\DeclareUnicodeCharacter{00D7}{\ensuremath{\times}}      % ×
\DeclareUnicodeCharacter{00E4}{\"a}             % ä
\DeclareUnicodeCharacter{00E9}{\'e}             % é
\DeclareUnicodeCharacter{00F6}{\"o}             % ö
\DeclareUnicodeCharacter{00FC}{\"u}             % ü
\DeclareUnicodeCharacter{010C}{\v{C}}           % Č
\DeclareUnicodeCharacter{010D}{\v{c}}           % č
\DeclareUnicodeCharacter{0394}{\ensuremath{\Delta}}      % Δ
\DeclareUnicodeCharacter{03A3}{\ensuremath{\Sigma}}      % Σ
\DeclareUnicodeCharacter{03A9}{\ensuremath{\Omega}}      % Ω
\DeclareUnicodeCharacter{03B1}{\ensuremath{\alpha}}      % α
\DeclareUnicodeCharacter{03B2}{\ensuremath{\beta}}       % β
\DeclareUnicodeCharacter{03B3}{\ensuremath{\gamma}}      % γ
\DeclareUnicodeCharacter{03B4}{\ensuremath{\delta}}      % δ
\DeclareUnicodeCharacter{03B5}{\ensuremath{\varepsilon}} % ε
\DeclareUnicodeCharacter{03BA}{\ensuremath{\kappa}}      % κ
\DeclareUnicodeCharacter{03BB}{\ensuremath{\lambda}}     % λ
\DeclareUnicodeCharacter{03BC}{\ensuremath{\mu}}         % μ
\DeclareUnicodeCharacter{03BD}{\ensuremath{\nu}}         % ν
\DeclareUnicodeCharacter{03C0}{\ensuremath{\pi}}         % π
\DeclareUnicodeCharacter{03C1}{\ensuremath{\rho}}        % ρ
\DeclareUnicodeCharacter{03C3}{\ensuremath{\sigma}}      % σ
\DeclareUnicodeCharacter{03C4}{\ensuremath{\tau}}        % τ
\DeclareUnicodeCharacter{03C6}{\ensuremath{\varphi}}     % φ
\DeclareUnicodeCharacter{03C7}{\ensuremath{\chi}}        % χ
\DeclareUnicodeCharacter{03C8}{\ensuremath{\psi}}        % ψ
\DeclareUnicodeCharacter{03C9}{\ensuremath{\omega}}      % ω
\DeclareUnicodeCharacter{2013}{--}              % – (en dash)
\DeclareUnicodeCharacter{2014}{---}             % — (em dash)
\DeclareUnicodeCharacter{2018}{`}               % ‘
\DeclareUnicodeCharacter{2019}{'}               % ’
\DeclareUnicodeCharacter{201C}{``}              % “
\DeclareUnicodeCharacter{201D}{''}              % ”
\DeclareUnicodeCharacter{2070}{\ensuremath{^{0}}}        % ⁰
\DeclareUnicodeCharacter{2071}{\ensuremath{^{i}}}        % ⁱ
\DeclareUnicodeCharacter{2122}{\textsuperscript{TM}}     % ™
\DeclareUnicodeCharacter{2126}{\ensuremath{\Omega}}      % Ω (ohm sign)
\DeclareUnicodeCharacter{210F}{\ensuremath{\hbar}}       % ℏ
\DeclareUnicodeCharacter{2115}{\ensuremath{\mathbb{N}}}  % ℕ
\DeclareUnicodeCharacter{2102}{\ensuremath{\mathbb{C}}}  % ℂ
\DeclareUnicodeCharacter{211D}{\ensuremath{\mathbb{R}}}  % ℝ
\DeclareUnicodeCharacter{2124}{\ensuremath{\mathbb{Z}}}  % ℤ
\DeclareUnicodeCharacter{211A}{\ensuremath{\mathbb{Q}}}  % ℚ
\DeclareUnicodeCharacter{2192}{\ensuremath{\to}}         % →
\DeclareUnicodeCharacter{21D2}{\ensuremath{\Rightarrow}} % ⇒
\DeclareUnicodeCharacter{2200}{\ensuremath{\forall}}     % ∀
\DeclareUnicodeCharacter{2203}{\ensuremath{\exists}}     % ∃
\DeclareUnicodeCharacter{2208}{\ensuremath{\in}}         % ∈
\DeclareUnicodeCharacter{2209}{\ensuremath{\notin}}      % ∉
\DeclareUnicodeCharacter{2227}{\ensuremath{\wedge}}      % ∧
\DeclareUnicodeCharacter{2228}{\ensuremath{\vee}}        % ∨
\DeclareUnicodeCharacter{2229}{\ensuremath{\cap}}        % ∩
\DeclareUnicodeCharacter{222A}{\ensuremath{\cup}}        % ∪
\DeclareUnicodeCharacter{2245}{\ensuremath{\cong}}       % ≅
\DeclareUnicodeCharacter{2248}{\ensuremath{\approx}}     % ≈
\DeclareUnicodeCharacter{2260}{\ensuremath{\neq}}        % ≠
\DeclareUnicodeCharacter{2264}{\ensuremath{\leq}}        % ≤
\DeclareUnicodeCharacter{2265}{\ensuremath{\geq}}        % ≥
\DeclareUnicodeCharacter{22C5}{\ensuremath{\cdot}}       % ⋅
\DeclareUnicodeCharacter{2295}{\ensuremath{\oplus}}      % ⊕
\DeclareUnicodeCharacter{2297}{\ensuremath{\otimes}}     % ⊗

% ---- Mathematics & symbols ------------------------------------------
\usepackage{amsmath,amssymb,bm,mathrsfs,mathtools}
\usepackage{amsthm}             % theorem environments (load BEFORE hyperref)

% ---- Theorem environments (shared counter set, TECT canonical) -------
% All theorems / lemmas / propositions / corollaries / remarks share a
% single counter so cross-references read consistently across a paper.
% Definitions get a separate counter (definitions are numbered in their
% own sequence by editorial convention).
\newtheorem{theorem}{Theorem}
\newtheorem{lemma}[theorem]{Lemma}
\newtheorem{proposition}[theorem]{Proposition}
\newtheorem{corollary}[theorem]{Corollary}

\theoremstyle{remark}
\newtheorem{remark}[theorem]{Remark}
\newtheorem{example}[theorem]{Example}

\theoremstyle{definition}
\newtheorem{definition}[theorem]{Definition}
\newtheorem{conjecture}[theorem]{Conjecture}
\newtheorem{hypothesis}[theorem]{Hypothesis}

% ---- Tables / floats / graphics --------------------------------------
\usepackage{booktabs}           % \toprule / \midrule / \bottomrule
\usepackage{array}
\usepackage{graphicx}
\usepackage{enumitem}           % \begin{itemize}[leftmargin=...] options
\usepackage{hyphenat}           % \hyp{} for breakable hyphens in \texttt
\usepackage{seqsplit}           % \seqsplit{} for long unbreakable strings
\sloppy                         % allow stretchier inter-word spacing to
                                % reduce overfull \hbox warnings on long URLs
                                % and long \texttt tokens (paragraph-level).
\emergencystretch=3em           % last-resort hbox stretch budget

% ---- Cross-references (hyperref last among the standard packages,
%      cleveref AFTER hyperref) ----------------------------------------
\usepackage[%
  colorlinks=true,%
  linkcolor=blue,%
  citecolor=blue,%
  urlcolor=blue,%
  breaklinks=true,%
  bookmarks=true,%
  bookmarksnumbered=true%
]{hyperref}
\usepackage[capitalise,nameinlink]{cleveref}

% ---- TECT-canonical author block macros -----------------------------
% (defined here so every paper has the same author / affiliation style)
\newcommand{\tectauthor}{%
  \author{Jusang Lee}%
  \email{jtkor@outlook.com}%
  \affiliation{Independent researcher, Republic of Korea}%
}

% ---- TECT-canonical archive footer ----------------------------------
\newcommand{\tectarchivefooter}{%
  \par\medskip
  \noindent\small\textit{Canonical archive}:
  \url{https://github.com/TECT-OpenPhysics/TECT}.\par%
}

% ---- TECT TODO / AUDIT-FLAG / HONEST-SCOPE inline marks --------------
% Used in drafts to highlight pending audit items without disturbing
% the body text style. Suppress via \tectsuppressmarks (define empty).
\providecommand{\tectauditflag}[1]{\textcolor{red}{\textbf{[AUDIT-FLAG: #1]}}}
\providecommand{\tecttodo}[1]{\textcolor{orange}{\textbf{[TODO: #1]}}}
\providecommand{\tecthonestscope}[1]{\textit{Honest scope: #1.}}

% =====================================================================
% End of TECT canonical preamble
% =====================================================================

%   \begin{document}
%   ...
%   \end{document}
%
% Build (every paper): `pdflatex Paper-NN && pdflatex Paper-NN`
% (the second pass resolves \ref{} / \cref{} forward references).
% A bash/PowerShell wrapper that automates this is at
% Docs/papers/papers/_shared/build-paper.sh / build-paper.ps1.
% =====================================================================

% ---- Document class --------------------------------------------------
% PRD class option of revtex4-2 with [reprint,onecolumn] gives the same
% APS heading / by-line / abstract style as PRL but in a wider single
% column that handles long display equations and theorem environments
% without column-break artefacts. nofootinbib keeps citations out of
% footnotes. The 11pt option restores standard font metrics; with T1
% fontenc + lmodern this gives non-bitmap glyphs.
\documentclass[%
  aps,prd,reprint,onecolumn,nofootinbib,11pt,%
  amsmath,amssymb,floatfix,longbibliography%
]{revtex4-2}

% ---- Encoding & fonts ------------------------------------------------
\usepackage[T1]{fontenc}        % proper 8-bit font encoding (no font corruption)
\usepackage[utf8]{inputenc}     % allow UTF-8 source
\usepackage{lmodern}            % scalable Latin Modern (replaces bitmap CM)
\usepackage{textcomp}           % extra text symbols
\usepackage{microtype}          % subtle kerning improvements

% ---- UTF-8 → LaTeX character mappings --------------------------------
% Maps the Unicode characters that occur in TECT body text (legacy from
% the chat-content auto-archival rule, CLAUDE.md §4) to LaTeX-safe
% commands. Without these declarations, T1+inputenc rejects the chars
% with "Unicode character X not set up for use with LaTeX".
% Adding declarations centrally (here) is preferable to per-file edits.
\DeclareUnicodeCharacter{00A7}{\S}              % §  (section sign)
\DeclareUnicodeCharacter{00B2}{\ensuremath{^{2}}}        % ²
\DeclareUnicodeCharacter{00B3}{\ensuremath{^{3}}}        % ³
\DeclareUnicodeCharacter{00B5}{\ensuremath{\mu}}         % µ (micro)
\DeclareUnicodeCharacter{00B7}{\ensuremath{\cdot}}       % · (middle dot)
\DeclareUnicodeCharacter{00D7}{\ensuremath{\times}}      % ×
\DeclareUnicodeCharacter{00E4}{\"a}             % ä
\DeclareUnicodeCharacter{00E9}{\'e}             % é
\DeclareUnicodeCharacter{00F6}{\"o}             % ö
\DeclareUnicodeCharacter{00FC}{\"u}             % ü
\DeclareUnicodeCharacter{010C}{\v{C}}           % Č
\DeclareUnicodeCharacter{010D}{\v{c}}           % č
\DeclareUnicodeCharacter{0394}{\ensuremath{\Delta}}      % Δ
\DeclareUnicodeCharacter{03A3}{\ensuremath{\Sigma}}      % Σ
\DeclareUnicodeCharacter{03A9}{\ensuremath{\Omega}}      % Ω
\DeclareUnicodeCharacter{03B1}{\ensuremath{\alpha}}      % α
\DeclareUnicodeCharacter{03B2}{\ensuremath{\beta}}       % β
\DeclareUnicodeCharacter{03B3}{\ensuremath{\gamma}}      % γ
\DeclareUnicodeCharacter{03B4}{\ensuremath{\delta}}      % δ
\DeclareUnicodeCharacter{03B5}{\ensuremath{\varepsilon}} % ε
\DeclareUnicodeCharacter{03BA}{\ensuremath{\kappa}}      % κ
\DeclareUnicodeCharacter{03BB}{\ensuremath{\lambda}}     % λ
\DeclareUnicodeCharacter{03BC}{\ensuremath{\mu}}         % μ
\DeclareUnicodeCharacter{03BD}{\ensuremath{\nu}}         % ν
\DeclareUnicodeCharacter{03C0}{\ensuremath{\pi}}         % π
\DeclareUnicodeCharacter{03C1}{\ensuremath{\rho}}        % ρ
\DeclareUnicodeCharacter{03C3}{\ensuremath{\sigma}}      % σ
\DeclareUnicodeCharacter{03C4}{\ensuremath{\tau}}        % τ
\DeclareUnicodeCharacter{03C6}{\ensuremath{\varphi}}     % φ
\DeclareUnicodeCharacter{03C7}{\ensuremath{\chi}}        % χ
\DeclareUnicodeCharacter{03C8}{\ensuremath{\psi}}        % ψ
\DeclareUnicodeCharacter{03C9}{\ensuremath{\omega}}      % ω
\DeclareUnicodeCharacter{2013}{--}              % – (en dash)
\DeclareUnicodeCharacter{2014}{---}             % — (em dash)
\DeclareUnicodeCharacter{2018}{`}               % ‘
\DeclareUnicodeCharacter{2019}{'}               % ’
\DeclareUnicodeCharacter{201C}{``}              % “
\DeclareUnicodeCharacter{201D}{''}              % ”
\DeclareUnicodeCharacter{2070}{\ensuremath{^{0}}}        % ⁰
\DeclareUnicodeCharacter{2071}{\ensuremath{^{i}}}        % ⁱ
\DeclareUnicodeCharacter{2122}{\textsuperscript{TM}}     % ™
\DeclareUnicodeCharacter{2126}{\ensuremath{\Omega}}      % Ω (ohm sign)
\DeclareUnicodeCharacter{210F}{\ensuremath{\hbar}}       % ℏ
\DeclareUnicodeCharacter{2115}{\ensuremath{\mathbb{N}}}  % ℕ
\DeclareUnicodeCharacter{2102}{\ensuremath{\mathbb{C}}}  % ℂ
\DeclareUnicodeCharacter{211D}{\ensuremath{\mathbb{R}}}  % ℝ
\DeclareUnicodeCharacter{2124}{\ensuremath{\mathbb{Z}}}  % ℤ
\DeclareUnicodeCharacter{211A}{\ensuremath{\mathbb{Q}}}  % ℚ
\DeclareUnicodeCharacter{2192}{\ensuremath{\to}}         % →
\DeclareUnicodeCharacter{21D2}{\ensuremath{\Rightarrow}} % ⇒
\DeclareUnicodeCharacter{2200}{\ensuremath{\forall}}     % ∀
\DeclareUnicodeCharacter{2203}{\ensuremath{\exists}}     % ∃
\DeclareUnicodeCharacter{2208}{\ensuremath{\in}}         % ∈
\DeclareUnicodeCharacter{2209}{\ensuremath{\notin}}      % ∉
\DeclareUnicodeCharacter{2227}{\ensuremath{\wedge}}      % ∧
\DeclareUnicodeCharacter{2228}{\ensuremath{\vee}}        % ∨
\DeclareUnicodeCharacter{2229}{\ensuremath{\cap}}        % ∩
\DeclareUnicodeCharacter{222A}{\ensuremath{\cup}}        % ∪
\DeclareUnicodeCharacter{2245}{\ensuremath{\cong}}       % ≅
\DeclareUnicodeCharacter{2248}{\ensuremath{\approx}}     % ≈
\DeclareUnicodeCharacter{2260}{\ensuremath{\neq}}        % ≠
\DeclareUnicodeCharacter{2264}{\ensuremath{\leq}}        % ≤
\DeclareUnicodeCharacter{2265}{\ensuremath{\geq}}        % ≥
\DeclareUnicodeCharacter{22C5}{\ensuremath{\cdot}}       % ⋅
\DeclareUnicodeCharacter{2295}{\ensuremath{\oplus}}      % ⊕
\DeclareUnicodeCharacter{2297}{\ensuremath{\otimes}}     % ⊗

% ---- Mathematics & symbols ------------------------------------------
\usepackage{amsmath,amssymb,bm,mathrsfs,mathtools}
\usepackage{amsthm}             % theorem environments (load BEFORE hyperref)

% ---- Theorem environments (shared counter set, TECT canonical) -------
% All theorems / lemmas / propositions / corollaries / remarks share a
% single counter so cross-references read consistently across a paper.
% Definitions get a separate counter (definitions are numbered in their
% own sequence by editorial convention).
\newtheorem{theorem}{Theorem}
\newtheorem{lemma}[theorem]{Lemma}
\newtheorem{proposition}[theorem]{Proposition}
\newtheorem{corollary}[theorem]{Corollary}

\theoremstyle{remark}
\newtheorem{remark}[theorem]{Remark}
\newtheorem{example}[theorem]{Example}

\theoremstyle{definition}
\newtheorem{definition}[theorem]{Definition}
\newtheorem{conjecture}[theorem]{Conjecture}
\newtheorem{hypothesis}[theorem]{Hypothesis}

% ---- Tables / floats / graphics --------------------------------------
\usepackage{booktabs}           % \toprule / \midrule / \bottomrule
\usepackage{array}
\usepackage{graphicx}
\usepackage{enumitem}           % \begin{itemize}[leftmargin=...] options
\usepackage{hyphenat}           % \hyp{} for breakable hyphens in \texttt
\usepackage{seqsplit}           % \seqsplit{} for long unbreakable strings
\sloppy                         % allow stretchier inter-word spacing to
                                % reduce overfull \hbox warnings on long URLs
                                % and long \texttt tokens (paragraph-level).
\emergencystretch=3em           % last-resort hbox stretch budget

% ---- Cross-references (hyperref last among the standard packages,
%      cleveref AFTER hyperref) ----------------------------------------
\usepackage[%
  colorlinks=true,%
  linkcolor=blue,%
  citecolor=blue,%
  urlcolor=blue,%
  breaklinks=true,%
  bookmarks=true,%
  bookmarksnumbered=true%
]{hyperref}
\usepackage[capitalise,nameinlink]{cleveref}

% ---- TECT-canonical author block macros -----------------------------
% (defined here so every paper has the same author / affiliation style)
\newcommand{\tectauthor}{%
  \author{Jusang Lee}%
  \email{jtkor@outlook.com}%
  \affiliation{Independent researcher, Republic of Korea}%
}

% ---- TECT-canonical archive footer ----------------------------------
\newcommand{\tectarchivefooter}{%
  \par\medskip
  \noindent\small\textit{Canonical archive}:
  \url{https://github.com/TECT-OpenPhysics/TECT}.\par%
}

% ---- TECT TODO / AUDIT-FLAG / HONEST-SCOPE inline marks --------------
% Used in drafts to highlight pending audit items without disturbing
% the body text style. Suppress via \tectsuppressmarks (define empty).
\providecommand{\tectauditflag}[1]{\textcolor{red}{\textbf{[AUDIT-FLAG: #1]}}}
\providecommand{\tecttodo}[1]{\textcolor{orange}{\textbf{[TODO: #1]}}}
\providecommand{\tecthonestscope}[1]{\textit{Honest scope: #1.}}

% =====================================================================
% End of TECT canonical preamble
% =====================================================================

%   \begin{document}
%   ...
%   \end{document}
%
% Build (every paper): `pdflatex Paper-NN && pdflatex Paper-NN`
% (the second pass resolves \ref{} / \cref{} forward references).
% A bash/PowerShell wrapper that automates this is at
% Docs/papers/papers/_shared/build-paper.sh / build-paper.ps1.
% =====================================================================

% ---- Document class --------------------------------------------------
% PRD class option of revtex4-2 with [reprint,onecolumn] gives the same
% APS heading / by-line / abstract style as PRL but in a wider single
% column that handles long display equations and theorem environments
% without column-break artefacts. nofootinbib keeps citations out of
% footnotes. The 11pt option restores standard font metrics; with T1
% fontenc + lmodern this gives non-bitmap glyphs.
\documentclass[%
  aps,prd,reprint,onecolumn,nofootinbib,11pt,%
  amsmath,amssymb,floatfix,longbibliography%
]{revtex4-2}

% ---- Encoding & fonts ------------------------------------------------
\usepackage[T1]{fontenc}        % proper 8-bit font encoding (no font corruption)
\usepackage[utf8]{inputenc}     % allow UTF-8 source
\usepackage{lmodern}            % scalable Latin Modern (replaces bitmap CM)
\usepackage{textcomp}           % extra text symbols
\usepackage{microtype}          % subtle kerning improvements

% ---- UTF-8 → LaTeX character mappings --------------------------------
% Maps the Unicode characters that occur in TECT body text (legacy from
% the chat-content auto-archival rule, CLAUDE.md §4) to LaTeX-safe
% commands. Without these declarations, T1+inputenc rejects the chars
% with "Unicode character X not set up for use with LaTeX".
% Adding declarations centrally (here) is preferable to per-file edits.
\DeclareUnicodeCharacter{00A7}{\S}              % §  (section sign)
\DeclareUnicodeCharacter{00B2}{\ensuremath{^{2}}}        % ²
\DeclareUnicodeCharacter{00B3}{\ensuremath{^{3}}}        % ³
\DeclareUnicodeCharacter{00B5}{\ensuremath{\mu}}         % µ (micro)
\DeclareUnicodeCharacter{00B7}{\ensuremath{\cdot}}       % · (middle dot)
\DeclareUnicodeCharacter{00D7}{\ensuremath{\times}}      % ×
\DeclareUnicodeCharacter{00E4}{\"a}             % ä
\DeclareUnicodeCharacter{00E9}{\'e}             % é
\DeclareUnicodeCharacter{00F6}{\"o}             % ö
\DeclareUnicodeCharacter{00FC}{\"u}             % ü
\DeclareUnicodeCharacter{010C}{\v{C}}           % Č
\DeclareUnicodeCharacter{010D}{\v{c}}           % č
\DeclareUnicodeCharacter{0394}{\ensuremath{\Delta}}      % Δ
\DeclareUnicodeCharacter{03A3}{\ensuremath{\Sigma}}      % Σ
\DeclareUnicodeCharacter{03A9}{\ensuremath{\Omega}}      % Ω
\DeclareUnicodeCharacter{03B1}{\ensuremath{\alpha}}      % α
\DeclareUnicodeCharacter{03B2}{\ensuremath{\beta}}       % β
\DeclareUnicodeCharacter{03B3}{\ensuremath{\gamma}}      % γ
\DeclareUnicodeCharacter{03B4}{\ensuremath{\delta}}      % δ
\DeclareUnicodeCharacter{03B5}{\ensuremath{\varepsilon}} % ε
\DeclareUnicodeCharacter{03BA}{\ensuremath{\kappa}}      % κ
\DeclareUnicodeCharacter{03BB}{\ensuremath{\lambda}}     % λ
\DeclareUnicodeCharacter{03BC}{\ensuremath{\mu}}         % μ
\DeclareUnicodeCharacter{03BD}{\ensuremath{\nu}}         % ν
\DeclareUnicodeCharacter{03C0}{\ensuremath{\pi}}         % π
\DeclareUnicodeCharacter{03C1}{\ensuremath{\rho}}        % ρ
\DeclareUnicodeCharacter{03C3}{\ensuremath{\sigma}}      % σ
\DeclareUnicodeCharacter{03C4}{\ensuremath{\tau}}        % τ
\DeclareUnicodeCharacter{03C6}{\ensuremath{\varphi}}     % φ
\DeclareUnicodeCharacter{03C7}{\ensuremath{\chi}}        % χ
\DeclareUnicodeCharacter{03C8}{\ensuremath{\psi}}        % ψ
\DeclareUnicodeCharacter{03C9}{\ensuremath{\omega}}      % ω
\DeclareUnicodeCharacter{2013}{--}              % – (en dash)
\DeclareUnicodeCharacter{2014}{---}             % — (em dash)
\DeclareUnicodeCharacter{2018}{`}               % ‘
\DeclareUnicodeCharacter{2019}{'}               % ’
\DeclareUnicodeCharacter{201C}{``}              % “
\DeclareUnicodeCharacter{201D}{''}              % ”
\DeclareUnicodeCharacter{2070}{\ensuremath{^{0}}}        % ⁰
\DeclareUnicodeCharacter{2071}{\ensuremath{^{i}}}        % ⁱ
\DeclareUnicodeCharacter{2122}{\textsuperscript{TM}}     % ™
\DeclareUnicodeCharacter{2126}{\ensuremath{\Omega}}      % Ω (ohm sign)
\DeclareUnicodeCharacter{210F}{\ensuremath{\hbar}}       % ℏ
\DeclareUnicodeCharacter{2115}{\ensuremath{\mathbb{N}}}  % ℕ
\DeclareUnicodeCharacter{2102}{\ensuremath{\mathbb{C}}}  % ℂ
\DeclareUnicodeCharacter{211D}{\ensuremath{\mathbb{R}}}  % ℝ
\DeclareUnicodeCharacter{2124}{\ensuremath{\mathbb{Z}}}  % ℤ
\DeclareUnicodeCharacter{211A}{\ensuremath{\mathbb{Q}}}  % ℚ
\DeclareUnicodeCharacter{2192}{\ensuremath{\to}}         % →
\DeclareUnicodeCharacter{21D2}{\ensuremath{\Rightarrow}} % ⇒
\DeclareUnicodeCharacter{2200}{\ensuremath{\forall}}     % ∀
\DeclareUnicodeCharacter{2203}{\ensuremath{\exists}}     % ∃
\DeclareUnicodeCharacter{2208}{\ensuremath{\in}}         % ∈
\DeclareUnicodeCharacter{2209}{\ensuremath{\notin}}      % ∉
\DeclareUnicodeCharacter{2227}{\ensuremath{\wedge}}      % ∧
\DeclareUnicodeCharacter{2228}{\ensuremath{\vee}}        % ∨
\DeclareUnicodeCharacter{2229}{\ensuremath{\cap}}        % ∩
\DeclareUnicodeCharacter{222A}{\ensuremath{\cup}}        % ∪
\DeclareUnicodeCharacter{2245}{\ensuremath{\cong}}       % ≅
\DeclareUnicodeCharacter{2248}{\ensuremath{\approx}}     % ≈
\DeclareUnicodeCharacter{2260}{\ensuremath{\neq}}        % ≠
\DeclareUnicodeCharacter{2264}{\ensuremath{\leq}}        % ≤
\DeclareUnicodeCharacter{2265}{\ensuremath{\geq}}        % ≥
\DeclareUnicodeCharacter{22C5}{\ensuremath{\cdot}}       % ⋅
\DeclareUnicodeCharacter{2295}{\ensuremath{\oplus}}      % ⊕
\DeclareUnicodeCharacter{2297}{\ensuremath{\otimes}}     % ⊗

% ---- Mathematics & symbols ------------------------------------------
\usepackage{amsmath,amssymb,bm,mathrsfs,mathtools}
\usepackage{amsthm}             % theorem environments (load BEFORE hyperref)

% ---- Theorem environments (shared counter set, TECT canonical) -------
% All theorems / lemmas / propositions / corollaries / remarks share a
% single counter so cross-references read consistently across a paper.
% Definitions get a separate counter (definitions are numbered in their
% own sequence by editorial convention).
\newtheorem{theorem}{Theorem}
\newtheorem{lemma}[theorem]{Lemma}
\newtheorem{proposition}[theorem]{Proposition}
\newtheorem{corollary}[theorem]{Corollary}

\theoremstyle{remark}
\newtheorem{remark}[theorem]{Remark}
\newtheorem{example}[theorem]{Example}

\theoremstyle{definition}
\newtheorem{definition}[theorem]{Definition}
\newtheorem{conjecture}[theorem]{Conjecture}
\newtheorem{hypothesis}[theorem]{Hypothesis}

% ---- Tables / floats / graphics --------------------------------------
\usepackage{booktabs}           % \toprule / \midrule / \bottomrule
\usepackage{array}
\usepackage{graphicx}
\usepackage{enumitem}           % \begin{itemize}[leftmargin=...] options
\usepackage{hyphenat}           % \hyp{} for breakable hyphens in \texttt
\usepackage{seqsplit}           % \seqsplit{} for long unbreakable strings
\sloppy                         % allow stretchier inter-word spacing to
                                % reduce overfull \hbox warnings on long URLs
                                % and long \texttt tokens (paragraph-level).
\emergencystretch=3em           % last-resort hbox stretch budget

% ---- Cross-references (hyperref last among the standard packages,
%      cleveref AFTER hyperref) ----------------------------------------
\usepackage[%
  colorlinks=true,%
  linkcolor=blue,%
  citecolor=blue,%
  urlcolor=blue,%
  breaklinks=true,%
  bookmarks=true,%
  bookmarksnumbered=true%
]{hyperref}
\usepackage[capitalise,nameinlink]{cleveref}

% ---- TECT-canonical author block macros -----------------------------
% (defined here so every paper has the same author / affiliation style)
\newcommand{\tectauthor}{%
  \author{Jusang Lee}%
  \email{jtkor@outlook.com}%
  \affiliation{Independent researcher, Republic of Korea}%
}

% ---- TECT-canonical archive footer ----------------------------------
\newcommand{\tectarchivefooter}{%
  \par\medskip
  \noindent\small\textit{Canonical archive}:
  \url{https://github.com/TECT-OpenPhysics/TECT}.\par%
}

% ---- TECT TODO / AUDIT-FLAG / HONEST-SCOPE inline marks --------------
% Used in drafts to highlight pending audit items without disturbing
% the body text style. Suppress via \tectsuppressmarks (define empty).
\providecommand{\tectauditflag}[1]{\textcolor{red}{\textbf{[AUDIT-FLAG: #1]}}}
\providecommand{\tecttodo}[1]{\textcolor{orange}{\textbf{[TODO: #1]}}}
\providecommand{\tecthonestscope}[1]{\textit{Honest scope: #1.}}

% =====================================================================
% End of TECT canonical preamble
% =====================================================================


\begin{document}

\title{BCC ranking and a numerical mass-gap anchor in a Brazovskii shell
model for Pillar 1 of topological energy condensate theory}

\author{Jusang Lee}
\email{jtkor@outlook.com}
\affiliation{Independent researcher, Republic of Korea}

\date{\today}

\begin{abstract}
We present a conditional Pillar-1 programme for the Topological Energy
Condensate Theory (TECT) combining (i) a BCC-favouring shell-model
ranking among standard 3D crystallographic competitors within the
single-mode-approximation (SMA) / equal-amplitude / first-shell
Brazovskii truncation~\cite{Math01v2}; (ii) a finite-$N$ numerical
mass-gap anchor at the Brazovskii operating point $\mu^2 = +5 \times
10^{-3}$ using a 32-mode spectral discretisation, yielding
$m^{*2} \approx 4.247 \times 10^{-2}$ in TECT code units (single
finite-$N$ run, NOT a continuum-limit certified value)~\cite{Math82AddF};
and (iii) an analytic shell-spectral-gap bound combined with
guarded-quotient acceptance criteria~\cite{Math82AddH}. The combination
constitutes a Pillar-1 \emph{programme}, NOT a completed closure
theorem: the BCC-vs-competitor ranking is conditional on the
single-shell SMA truncation; the numerical anchor is a finite-$N$
single-operating-point measurement (not a global-minimum proof); and
the analytic bound is a shell-gap estimate, not a unique-minimum
theorem. The current canonical tier of Pillar 1 is recorded in
\texttt{Docs/status/TOE-FACT-SHEET.md}; a pre-registered F1.1
falsification gate $|\Delta m^{*2} / m^{*2}_{\rm analytic}| < 0.1$
binds the numerical anchor against the analytic estimate.
\end{abstract}

\maketitle

% =====================================================================
\section{Introduction}
% =====================================================================
The Theory-of-Everything (TOE) candidate framework known as Topological
Energy Condensate Theory (TECT) posits an 11-pillar structural
programme~\cite{Math60}. Pillar~1 --- the BCC ground-state ranking
and mass-gap extraction --- has historically been treated as a
scaffold-tier result based on qualitative resonance arguments from
condensed-matter phase transitions. The present work organises three
mutually consistent components into a single conditional Pillar-1
\emph{programme}: (i) a single-shell-SMA / equal-amplitude
free-energy ranking among standard 3D crystallographic competitors,
(ii) an explicit finite-$N$ numerical mass-gap anchor at the TECT
operating point, and (iii) an analytic shell-spectral-gap bound. The
combination is a Pillar-1 programme rather than a closure theorem;
the canonical Pillar-1 tier is recorded in
\texttt{Docs/status/TOE-FACT-SHEET.md}.

The Brazovskii model describes the isotropic-to-modulated phase
transition in systems with competing interactions~\cite{Math97}. TECT
employs this framework to describe the candidate primordial BCC
condensate from which both spacetime geometry and quantum mechanics
are conjectured to emerge. Establishing that BCC is the lowest-ranked
member of the tested competitor set within the single-shell SMA
truncation is the present scope; certifying that BCC is the global
minimum of the unrestricted Brazovskii configuration space is a
strictly larger open programme that would additionally require
beyond-mean-field control and beyond-single-shell analysis. The
present paper does not address that larger programme; it addresses
the conditional ranking, the numerical anchor, and the analytic
bound, each of which carries explicit hypotheses stated below.

% =====================================================================
\section{Setting and main result}
% =====================================================================
\subsection{Brazovskii functional and locked parameters}

The TECT condensate order parameter $\Psi : \mathbb{R}^3 \to \mathbb{C}^n$
obeys the Brazovskii free energy
\begin{equation}
\mathcal{F}[\Psi] = \int d^3x \left[
\frac{1}{2}|\nabla\Psi|^2 + \frac{1}{2}\mu^2|\Psi|^2 +
\frac{\gamma}{2}(\Delta + q_0^2)^2|\Psi|^2 + \frac{\lambda}{4}|\Psi|^4
\right],
\label{eq:brazovskii}
\end{equation}
with parameters locked by the TECT operating-point gate Math195:
$(\mu^2, \lambda, \gamma) = (+5\times 10^{-3}, -0.43, 1.62)$ and shell
radius $q_0 = 0.6802$.

The TECT framework postulates that this functional governs the primordial
universe from its pre-condensation symmetric phase through Kibble--Zurek
freeze-out to the locked BCC state, and that all Standard Model degrees of
freedom emerge from topological defects in this condensate.

\subsection{Main result (conditional)}

\begin{theorem}[Conditional Pillar-1 ranking and numerical mass-gap
anchor]
\label{thm:main}
Let $\mathcal{F}[\Psi]$ be the Brazovskii functional
Eq.~\eqref{eq:brazovskii} at the TECT operating point. Within the
single-shell single-mode-approximation (SMA), equal-amplitude,
first-shell truncation, restricted to the nine tested standard 3D
crystallographic ans\"atze (lamellar, square columnar, hexagonal
columnar, simple cubic, FCC, HCP, A15 Frank--Kasper, gyroid,
$\sigma$-phase), the BCC ans\"atz attains the lowest mean-field
free-energy density $\mathcal{F}[\Psi]/V$ among the tested set.

At grid resolution $N = 32$ per lattice constant, the numerical
Newton--Krylov continuation produces a finite-$N$ candidate broken-phase
configuration whose projected (transverse) Hessian, after removing the
three translational zero modes, exhibits a positive lowest eigenvalue
\[
\lambda_{\min}(\widetilde{\mathrm{Hess}})
= 4.247 \times 10^{-2}\quad\text{in TECT code units},
\]
which we adopt as the finite-$N$ mass-gap \emph{anchor}
\[
m^{*2}_{\rm anchor}(N=32) = 4.247 \times 10^{-2}.
\]
This finite-$N$ anchor is consistent with the analytic
shell-spectral-gap estimate $m^{*2}_{\rm analytic} \approx 9.005$
within the H1--H3 hypotheses below; the residual $\sim 200\times$
ratio reflects the conservative bound character of the analytic
estimate (which is a shell-gap upper bound rather than an exact
identity), not a precision claim.

\textbf{Honest scope.} The ranking is conditional on the SMA /
equal-amplitude / first-shell truncation; the numerical anchor is a
single-operating-point finite-$N$ measurement (not a continuum-limit
certified value); the analytic estimate is a shell-spectral-gap upper
bound (not a unique-minimum theorem). The combined statement is a
Pillar-1 \emph{programme}, not a closure theorem on the unrestricted
Brazovskii configuration space.
\end{theorem}

\begin{remark}[On the Hessian positivity claim]
\label{rem:hessian}
The positive-definiteness of the projected Hessian at $N=32$ is an
empirical finding from the Newton continuation, not a theorem on the
continuum Brazovskii functional. In particular, it is conditional on
(i) the choice of zero-mode projector (here only the three
translation modes; rotational and internal-phase modes are not
projected out), and (ii) the finite resolution $N=32$. The Math292
4-gauge acceptance criterion (in particular gauge G3, transverse
Hessian PSD on the full soft-mode complement) is not yet certified by
the present finite-$N$ anchor and is tracked separately as part of
the F-Pillar6 verdict programme (deadline 2026-05-29).
\end{remark}

\section{Three-component framework (after Math82-AddH)}\label{sec:three-component}
% =====================================================================

The Pillar-1 conditional programme synthesises three independently
verified components per the Math82-AddH structure:

\paragraph{Component 1 — Single-shell SMA ranking with multi-shell
robustness (Math01-v2 + companion paper TI-2).}
\cite{Math01v2} (BCC uniqueness rigorous) systematically enumerates
the nine standard 3D crystallographic competitors under the
single-mode approximation (SMA) with fixed shell radius $q_0$. For
each competitor, the mean-field free-energy density
$\mathcal{F}_{\rm MF}[\{A_i\}]/V$ is computed; at the
equal-amplitude point (protected by Michel's theorem), the
Brazovskii cubic-stabilisation coefficient
$\mathcal{C}_{\rm lat} = N_{\rm loop}/N_S^2$ determines the
phase hierarchy. The exact combinatorial counts (companion top-impact
paper TI-2 Theorem on cubic Bravais lattices) give $\mathcal{C}_{\rm SC}
= 0$, $\mathcal{C}_{\rm FCC} = 1/8$, $\mathcal{C}_{\rm BCC} = 1/6$.
The single-shell ranking is upgraded to multi-shell robustness in the
strong primary-shell-dominance regime $\kappa = |r_1|/|r_0| \gg 1$
(TI-2 Multi-shell BCC Persistence Theorem) and to global 12-star
optimality $\Lambda(S_{\rm BCC}) = 540$ on antipodal 12-vector
configurations of $\mathbb{S}^2$ (TI-2 Global 12-Star Optimality
Theorem).

\paragraph{Component 2 — Finite-$N$ numerical mass-gap anchor
(Math82-AddF / Math82-AddH-Result-EXECUTED).}
\cite{Math82AddF} executes the spectral PDE continuation from
$\mu^2 = -1$ to $\mu^2 = +5 \times 10^{-3}$ on the
maximally-symmetric 6-cosine BCC seed. With extended Newton budget
(20 iterations), the continuation reaches the convergence-residual
$\|\nabla\mathcal{F}\|/\sqrt{\dim} = 6.11 \times 10^{-9}$ at
$N = 32$. The Phase-2 Lanczos iteration on the projected Hessian
(translation zero-modes removed) reports $\lambda_{\min}^{(\perp)}
= 4.247 \times 10^{-2}$ at the converged iterate. We adopt this as
the finite-$N$ anchor
\[
m^{*2}_{\rm anchor}(N=32) = 4.247 \times 10^{-2} \quad \text{(TECT
code units)},
\]
with the conditional-character caveat of
Remark~\ref{rem:hessian}.

\paragraph{Component 3 — Analytic shell-spectral-gap bound
(Math82-AddH analytical derivation).}
\cite{Math82AddH} derives an analytical bound on $m^{*2}$ from the
Brazovskii shell-spectral-gap structure. The shell-mode spectrum
decomposes into three regimes: IR modes ($|\mathbf{k}| < \delta$,
condensate order parameter), shell modes ($|\mathbf{k} - q_0| <
\Delta_s$, Goldstone--Nambu excitations breaking $SO(3) \to O(2)$),
and UV modes ($|\mathbf{k}| > q_0 + \delta$, gapped at the critical
point). The shell modes define a spectral gap
\[
\Delta_{\rm shell} \sim |\mu^2| + O(h^2), \quad h = a_{\rm BCC}/N,
\]
separating IR from UV. The Math82-AddH derivation, combining the
shell-spectral-gap structure with the Class-II guarded-quotient
acceptance gate $\rho_{\rm cut} \sim 10^{-3}\, \varphi_0^2$
(Math56-AddB \S 2.6), yields the analytic estimate
\[
m^{*2}_{\rm analytic} = 9.005 \quad \text{(TECT code units)}.
\]
This is a shell-gap upper bound, NOT a unique-minimum theorem; the
factor of $\sim 200$ between the analytic upper bound and the
finite-$N$ anchor reflects the conservative-bound character of the
analytic estimate (the bound is on a quantity larger than the
actual gap), not a precision discrepancy.

\paragraph{Math82-AddH conditional-tier closure (three load-bearing
hypotheses).}
The Math82-AddH PROVED CONDITIONAL claim depends on three
explicitly enumerated hypotheses:
\begin{itemize}
\item \textbf{H1} (Brazovskii shell-spectral-gap analysis valid on
the BCC lattice for $N \geq 32$ with $|q_0 h| \lesssim 0.1$);
\item \textbf{H2} (the v2.4 continuation endpoint is a global
minimum: $\|\nabla\mathcal{F}\| < 10^{-8}$ + projected
$\lambda_{\min} > 0$ + shell-spectral-gap structure as predicted by
H1);
\item \textbf{H3} (Phase-2.5 Gates G0--G3 implementable: G0
$\|\Psi^*\|_{\rm RMS}/\varphi_0$ RMS gate, G1 spectral positivity
via Lanczos, G2 Ritz-vector overlap, G3 guarded-quotient,
all $N$-independent).
\end{itemize}
Under H1--H3, the three components combine to support the
conditional Pillar-1 programme.

\paragraph{Math97 universality-class anchor.}
The Math97 conditional Brazovskii-class membership programme
(companion auxiliary paper Auxiliary-01) provides the candidate
anchor: under the Brazovskii axioms (C1)--(C5) and the
obstruction bounds (O1)--(O5), TECT belongs to the
$\mathcal{U}_{\rm iso\text{-}Brazovskii}$ universality class. The
Math82-AddH global-minimum guarantee in the admissible regime
$\mu^2 \in [-1, 0.05]$ rests on this universality-class membership;
its conditional character is the Math97 conditional-programme
caveat (treated in Auxiliary-01 in detail).

\paragraph{Step 4 (Hypothesis validation).}
The proof is conditional on hypotheses: (H1) Brazovskii shell-spectral-gap
analysis holds on the BCC lattice for $N \geq 32$; (H2) the continuation
endpoint is a true global minimum; (H3) Phase-2.5 gates are implementable
to specified tolerances.

% =====================================================================
\section{Numerical verification}
% =====================================================================
\subsection{Run summary}

The definitive numerical run \cite{Math82AddF} executed the command:
\begin{verbatim}
python Codes\pde\continuation_mu2_v25.py
  --N 32  --L 62.20036  --mu2 5e-3
  --max-newton 20  --tol-newton 1e-8
\end{verbatim}
on the TECT private compute cluster. The run converged in 14 Newton iterations
to a stable minimum.

\subsection{Key extracted quantities}

\begin{table}[h]
\centering
\begin{tabular}{llr}
\hline
Quantity & Description & Value \\
\hline
$\|\nabla\mathcal{F}\|/\sqrt{\dim}$ & gradient norm (convergence check) & $6.11 \times 10^{-9}$ \\
$\lambda_{\min}$ & smallest Hessian eigenvalue & $+4.247 \times 10^{-2}$ \\
$m^{*2}$ & squared mass gap (main result) & $4.247 \times 10^{-2}$ \\
$\Delta F_{\rm num}$ & finite-$N$ vacuum preference & $+4.15 \times 10^{-10}$ \\
\hline
\end{tabular}
\end{table}

The small positive vacuum preference ($\Delta F_{\rm num} > 0$) is expected at
finite $N = 32$ and is flagged for verification at $N = 64$.

\subsection{Falsification gate}

F1.1: $|\Delta m^{*2}/m^{*2}_{\rm analytic}| < 0.1$ (deadline 2026-06-15).
Current status: OPEN.

% =====================================================================
\section{Discussion and outlook}
% =====================================================================
The conditional Pillar-1 programme established in
Theorem~\ref{thm:main} feeds into the subsequent pillars of the TECT
framework as a programme-level numerical anchor: Pillar~2 (kinematic
Lorentz consistency) builds on the BCC dispersion at the operating
point; Pillar~3 (emergent-gravity programme) uses the $m^{*2}$ scale
as a candidate normalisation; Pillars~4--7 (gauge sector) inherit
the BCC defect-bundle structure. The mass-gap anchor
$m^{*2}_{\rm anchor}(N=32) = 4.247 \times 10^{-2}$ is provided as a
finite-$N$ numerical reference, not as a continuum-limit certified
value.

\textbf{Honest scope.} The scope limitations of the present
programme, which are intrinsic and not removable by tightening the
present proof, are: (1) single-shell SMA truncation in the
competitor enumeration, which excludes multi-shell or
non-Bravais-lattice configurations; (2) finite-$N$ discretisation
effects, which the planned $N \in \{16, 32, 64\}$ ladder is
designed to address via Richardson extrapolation; (3) the
equal-amplitude ans\"atz at Michel-theorem-protected points, which
may miss asymmetric perturbations; (4) the conditional character of
the projected-Hessian positivity claim (Remark~\ref{rem:hessian}),
which is not yet a Math292 G3 acceptance certificate. Each of these
is an explicit open task in the canonical Pillar-1 bookkeeping; the
present programme does NOT claim to close them. The canonical tier
of Pillar~1 is recorded in
\texttt{Docs/status/TOE-FACT-SHEET.md} and remains conditional
pending the F1.1 falsification verdict and the planned multi-$N$
ladder closure.

\textbf{Implications for the foundational ``why BCC?'' question.}
At the level of the multi-mode-robust + globally-12-star-optimal
single-shell Brazovskii framework developed in the companion
top-impact paper TI-2 (which incorporates the multi-mode robustness,
exact combinatorial counting, global 12-star optimality, and
multi-shell BCC persistence theorems from the historical TECT-BCC-Part-I
draft), the BCC selection is rigorous in the combinatorial and
global-12-star senses, robust against small-detuning multi-mode
corrections, and persistent in the strong-shell-dominance regime.
The present Pillar-1 programme is thereby anchored on a four-layer
structural-selection certificate at the conditional-programme level,
without claim of an unrestricted-configuration-space global-minimum
theorem.

% =====================================================================
\begin{acknowledgments}
The author thanks the Anthropic Claude Opus 4.7 collaborator for rigorous
audit-discipline enforcement during the Pillar 1 closure programme. This work
is part of TECT, canonical archive at
\url{https://github.com/TECT-OpenPhysics/TECT}.
\end{acknowledgments}

% =====================================================================
\bibliographystyle{apsrev4-2}
\bibliography{../_shared/tect-references}
% =====================================================================

\end{document}
